\documentclass[11pt]{article}
\usepackage[utf8]{inputenc}
\usepackage{amsmath,amssymb}
\usepackage[margin=1in]{geometry}
\usepackage{hyperref}
\emergencystretch=3em

\title{Building $D_{-\nu}$: the negative-order parabolic cylinder function in Lean}
\author{Claude, with Daniel Murfet}
\date{2026-09-06}

\begin{document}
\maketitle
\sloppy

\begin{abstract}
Unit 6 of the grammar-paper \S4 autoformalisation, and the first step of the option-B plan to build
the parabolic-cylinder infrastructure Mathlib lacks. We define the negative-order function
$D_{-\nu}(x)=e^{-x^2/4}/\Gamma(\nu)\cdot\int_0^\infty u^{\nu-1}e^{-u^2/2-xu}\,du$ ($\nu>0$) via its
DLMF \S12.5(i) integral representation --- exactly the orders $D_{-2\lambda}$ that
\texttt{cor:fluctuation\_closed\_form} needs --- prove the DLMF integrand integrable on $(0,\infty)$,
and pin the value $D_{-\nu}(0)=2^{\nu/2-1}\Gamma(\nu/2)/\Gamma(\nu)$ as a sanity anchor. Zero
\texttt{sorry}/\texttt{axiom}.
\end{abstract}

\section{Setting}
Units 1--5 formalised the whole of \S4.1: Watanabe's fluctuation function $S_\lambda(a)$, its
integrability and base value, the derivative ladder, the IBP recurrence and ODE, the boundary case,
and the lowering identity. What remains of \S4 --- the closed form $S_\lambda(a)=2^{1-\lambda}
\beta^{-\lambda}\Gamma(2\lambda)e^{\beta a^2/8}D_{-2\lambda}(-a\sqrt{\beta/2})$ and the QHO
identification --- rests on the parabolic cylinder function $D_\nu$, which Mathlib does not have. So
the arc forks: rather than defer, we build the piece of $D_\nu$ theory the paper actually uses. The
negative-order integral representation is the right entry point --- it is elementary (a Gamma-type
integral, no Weber ODE or Hermite expansion needed) and it is precisely the branch the closed form
lands in, since $2\lambda>0$.

\section{The things formalised}
{\small
\begin{verbatim}
noncomputable def parCylNeg (ν x : ℝ) : ℝ :=
  Real.exp (-x ^ 2 / 4) / Real.Gamma ν *
    ∫ u in Set.Ioi (0 : ℝ), u ^ (ν - 1) * Real.exp (-u ^ 2 / 2 - x * u)

theorem parCylNeg_integrableOn (ν x : ℝ) (hν : 0 < ν) :
    IntegrableOn (fun u => u ^ (ν - 1) * Real.exp (-u ^ 2 / 2 - x * u)) (Set.Ioi 0)

theorem parCylNeg_zero (ν : ℝ) (hν : 0 < ν) :
    parCylNeg ν 0 = 2 ^ (ν / 2 - 1) * Real.Gamma (ν / 2) / Real.Gamma ν
\end{verbatim}
}

\section{Proof strategy}
\textbf{Integrability.} The integrand $u^{\nu-1}e^{-u^2/2-xu}$ is dominated by a pure Gaussian: the
phase satisfies $-u^2/2-xu\le x^2-u^2/4$, which is $(u/2+x)^2\ge0$ rearranged. So the integrand is
$\le e^{x^2}\cdot u^{\nu-1}e^{-u^2/4}$, and the majorant is integrable by
\texttt{integrableOn\_rpow\_mul\_exp\_neg\_mul\_rpow} (the general Gamma-integral integrability at
$p=2$, $q=\nu-1>-1$, $b=\tfrac14$). \texttt{Integrable.mono'} closes it, with the pointwise bound via
\texttt{nlinarith [sq\_nonneg (u/2+x)]} and the $\mathrm{rpow}$/$\mathrm{npow}$ bridge
$u^{(2:\mathbb R)}=u^2$.

\textbf{Value at $0$.} At $x=0$ the integral is the standard $\int_0^\infty u^{\nu-1}e^{-u^2/2}\,du
=2^{\nu/2-1}\Gamma(\nu/2)$, an instance of \texttt{integral\_rpow\_mul\_exp\_neg\_mul\_rpow} at
$p=2$, $b=\tfrac12$, giving $(\tfrac12)^{-\nu/2}\cdot\tfrac12\cdot\Gamma(\nu/2)$; the $\mathrm{rpow}$
identity $(\tfrac12)^{-\nu/2}=2^{\nu/2}$ and $2^{\nu/2-1}=2^{\nu/2}\cdot\tfrac12$ reconcile it to the
stated form. The prefactor $e^{-0^2/4}/\Gamma(\nu)=1/\Gamma(\nu)$ then gives the quotient.

\section{Roadblocks and resolutions}
All Lean bookkeeping, no mathematical friction:
\begin{itemize}
\item \texttt{integral\_rpow\_mul\_exp\_neg\_mul\_rpow} takes its hypotheses in the order $(0<p)$,
$(-1<q)$, $(0<b)$; feeding $q$'s bound where $p$'s was expected left a stray $\vdash 0<\nu$. The
\texttt{integrableOn} sibling happens to accept $q$'s bound first, so the two are \emph{not}
argument-compatible --- check each signature rather than copying.
\item The $\mathrm{rpow}$ arithmetic $(\tfrac12)^{-\nu/2}=2^{\nu/2}$ needs
\texttt{Real.inv\_rpow} then \texttt{Real.rpow\_neg}$+$\texttt{inv\_inv}; \texttt{Real.rpow\_neg\_one}
does not exist on this pin (a known repo gotcha), so $2^{-1}$ went through
\texttt{Real.rpow\_neg}$+$\texttt{Real.rpow\_one}.
\item \texttt{setIntegral\_congr\_fun} hands the pointwise goal as applied lambdas; \texttt{beta\_reduce}
before the \texttt{rw}, and rewrite the exp-argument directly (\texttt{ring} cannot pass under
\texttt{exp}).
\item The final $\mathrm{rw}$-chain closes the intermediate integral value by \texttt{rfl}, so a
trailing \texttt{ring} errored ``no goals'' --- deleted; the outer prefactor step keeps its
\texttt{ring}.
\end{itemize}

\section{What was Mathlib, what was new}
Both Gamma-integral engines (\texttt{integral\_}/\texttt{integrableOn\_rpow\_mul\_exp\_neg\_mul\_rpow})
are Mathlib. The new content is the definition itself --- Mathlib has no parabolic cylinder function ---
and the two elementary facts that make it usable downstream: an integrability certificate valid for
\emph{every} $x$ (via the completed-square Gaussian majorant, so later change-of-variables arguments
have their dominating function ready) and the closed value at the origin.

\section{Lessons}
When a paper needs a special function absent from the library, the negative-order / integral-representation
branch is usually the cheapest foothold: it is a Gamma-type integral, so the existing
$\int x^q e^{-bx^p}$ machinery does all the analytic work, and one completed square gives an
$x$-uniform majorant that unlocks the substitution proofs to come. Build the representation the target
formula actually consumes, not the ``canonical'' definition.

\section{Follow-ups}
Next unit: \texttt{cor:fluctuation\_closed\_form} itself --- the change of variables $u=\sqrt{2\beta t}$
in $S_\lambda$, matching the kernel against \texttt{parCylNeg}\,$(2\lambda)\,(-a\sqrt{\beta/2})$ (scale
by \texttt{integral\_comp\_mul\_left\_Ioi}, then \texttt{integral\_comp\_rpow\_Ioi\_of\_pos} backwards
at $p=2$, with Gamma left outside the substitution). Then \texttt{eq:Dn\_hermite} ($D_n$ at integer
order via Mathlib's Hermite polynomials) and \texttt{lem:qho}.
\end{document}
